\chapter{Discuss\~ao}
\label{cap:discussao}

Este cap\'itulo interpreta os resultados apresentados no Cap\'itulo~\ref{cap:resultados} \`a luz dos fundamentos te\'oricos do Cap\'itulo~\ref{cap:fundamentacao} e dos estudos relacionados do Cap\'itulo~\ref{cap:trabalhos}, respondendo \`a quest\~ao central do trabalho e situando as constata\c c\~oes emp\'iricas no contexto mais amplo da pesquisa em aprendizado de m\'aquina aplicado \`a cardiologia.

\section{A Quest\~ao Central Respondida}
\label{sec:questao_central}

A pergunta que orienta este trabalho, \textit{para dados cl\'inicos tabulares de cardiopatas, o MLP produz resultados expressivamente superiores ao ADALINE?}, pode ser respondida de forma objetiva com base nos 24 experimentos conduzidos: n\~ao, ao menos n\~ao de forma consistente e expressiva.

No Banco 1, a diferen\c ca entre o melhor ADALINE (NumPy C3: 86,96\%) e o melhor MLP em termos de acur\'acia (Sklearn C2: 88,04\%) \'e de apenas 1,08 pontos percentuais, margem que, para um conjunto de 276 pacientes de teste, corresponde \`a classifica\c c\~ao correta de apenas 3 pacientes adicionais. Em Falsos Negativos, a vantagem do MLP NumPy C3 sobre o ADALINE NumPy C3 \'e de 5 FNs (15 versus 20), o que \'e clinicamente mais relevante, mas tamb\'em representa uma diferen\c ca poss\'ivel de ser atribu\'ida \`a variabilidade estat\'istica de um conjunto de teste de tamanho moderado.

No Banco 2, o MLP n\~ao demonstrou superioridade consistente: em duas das tr\^es configura\c c\~oes (C2 e C3), o MLP Sklearn produziu resultados inferiores ao ADALINE em acur\'acia, e o MLP NumPy em C3 (71,11\%) ficou atr\'as do ADALINE NumPy C3 (75,56\%). Apenas na configura\c c\~ao C1, o MLP NumPy ofereceu o melhor equil\'ibrio global (78,89\%, 9 FNs), mas o modelo degradou nas demais configura\c c\~oes enquanto o ADALINE manteve progress\~ao estável ou estacion\'aria.

\section{Interpreta\c c\~ao Te\'orica dos Resultados}
\label{sec:interpretacao_teorica}

Os padr\~oes de desempenho observados nos 24 experimentos n\~ao emergem de forma aleat\'oria: eles s\~ao consequ\^encias diretas das propriedades matem\'aticas de cada arquitetura, do otimizador utilizado e do volume de dados dispon\'ivel. Esta se\c c\~ao interpreta os principais achados emp\'iricos do Cap\'itulo~\ref{cap:resultados} \`a luz dos fundamentos te\'oricos apresentados no Cap\'itulo~\ref{cap:fundamentacao}, articulando a teoria com a evid\^encia experimental.

\subsection{O ADALINE e a Convexidade da Fun\c c\~ao de Custo}
\label{sec:discussao_adaline}

A estabilidade e previsibilidade do ADALINE ao longo de todas as configura\c c\~oes testadas, sem qualquer caso de degrada\c c\~ao, decorrem diretamente de uma propriedade matem\'atica fundamental: a convexidade de sua fun\c c\~ao de custo. O Erro Quadr\'atico M\'edio de um neurônio com fun\c c\~ao de ativa\c c\~ao sigmoide n\~ao \'e estritamente convexo como no caso do modelo puramente linear, mas a regi\~ao de opera\c c\~ao relevante para classifica\c c\~ao bin\'aria \'e bem comportada e livre de m\'inimos locais significativos \cite{braga2007,silva2016}. Isso significa que o gradiente descendente, independentemente do ponto de inicializa\c c\~ao, converge para uma solu\c c\~ao est\'avel, e que diferentes configura\c c\~oes de hiperpar\^ametros alteram apenas a velocidade de converg\^encia, n\~ao o destino final.

Essa propriedade tem implica\c c\~oes pr\'aticas diretas: um modelo com fun\c c\~ao de custo bem comportada exige menos ajuste de hiperpar\^ametros e oferece maior previsibilidade em ambientes de uso real, onde retreinamentos peri\'odicos com novos dados s\~ao frequentes \cite{silva2016}. Nos experimentos deste trabalho, o ADALINE NumPy nunca produziu surpresas negativas, cada configura\c c\~ao mais longa resultou em desempenho igual ou melhor, o que o torna prefer\'ivel em cen\'arios onde a confiabilidade do comportamento \'e prioridade sobre a maximiza\c c\~ao da acur\'acia absoluta.

\subsection{O MLP e o Dilema da Capacidade Subutilizada}
\label{sec:discussao_mlp}

O Teorema da Aproxima\c c\~ao Universal garante que um MLP com uma camada oculta e fun\c c\~ao de ativa\c c\~ao n\~ao linear pode aproximar qualquer fun\c c\~ao cont\'inua em dom\'inios compactos com precis\~ao arbitr\'aria \cite{braga2007}. Entretanto, esse teorema \'e um resultado de exist\^encia, n\~ao de constru\c c\~ao: ele garante que existe uma configura\c c\~ao de pesos capaz de realizar a aproxima\c c\~ao desejada, mas n\~ao que o algoritmo de treinamento encontrar\'a essa configura\c c\~ao a partir de dados finitos; crucialmente, n\~ao especifica o volume de dados necess\'ario para aprender essa configura\c c\~ao com garantias de generaliza\c c\~ao.

Nos experimentos deste trabalho, o MLP NumPy no Banco 1 evidenciou que essa capacidade latente existe: a progress\~ao de 82,97\% em C1 para 87,68\% em C3, com redu\c c\~ao de 27 para 15 FNs, demonstra que o modelo estava genuinamente aprendendo representa\c c\~oes mais ricas ao longo do treinamento. Contudo, o mesmo modelo no Banco 2 mostrou que a capacidade superior se converte rapidamente em \textit{overfitting} quando o volume de dados \'e insuficiente para regularizar o aprendizado \cite{silva2016}. A tens\~ao entre capacidade de representa\c c\~ao e generaliza\c c\~ao, denominada dilema vi\'es-vari\^ancia na literatura de aprendizado de m\'aquina, fica claramente evidenciada nessa compara\c c\~ao.

O experimento complementar com dados brutos (Se\c c\~ao~\ref{sec:resultados_bruto}) acrescenta uma dimens\~ao interpretativa relevante a esse dilema. Mesmo com representa\c c\~ao cont\'inua das vari\'aveis, o ADALINE superou o MLP no Banco~2 em todas as configura\c c\~oes e apresentou desempenho muito pr\'oximo no Banco~1. Esse padr\~ao \'e consistente com a hip\'otese de que a fronteira de decis\~ao \'otima nesses dados \'e \textit{predominantemente linear}: se existissem estruturas n\~ao lineares expressivas e explor\'aveis, o ADALINE, incapaz de represent\'a-las por constru\c c\~ao, seria sistematicamente superado pelo MLP mesmo em regimes de baixo volume de dados. A competitividade persistente do ADALINE com dados cont\'inuos sugere que os padr\~oes cl\'inicos determinantes para a classifica\c c\~ao podem ser capturados por combina\c c\~oes lineares das vari\'aveis dispon\'iveis. N\~ao \'e poss\'ivel, contudo, descartar a hip\'otese alternativa: o volume insuficiente do Banco~2 pode estar impedindo o MLP de materializar qualquer vantagem n\~ao linear latente, tornando as duas explica\c c\~oes empiricamente indistingu\'iveis com os dados atuais \cite{silva2016,braga2007}.

A inicializa\c c\~ao de Xavier, utilizada no MLP NumPy, desempenhou papel relevante na qualidade do treinamento. Ao calibrar a vari\^ancia dos pesos iniciais em fun\c c\~ao do n\'umero de neur\^onios de entrada e sa\'ida ($\pm\sqrt{6/(n_{\text{ent}} + n_{\text{saí}})}$), essa estrat\'egia evita tanto a satura\c c\~ao das fun\c c\~oes de ativa\c c\~ao (pesos muito grandes) quanto o desaparecimento do gradiente (pesos muito pr\'oximos de zero) nas fases iniciais do treinamento \cite{braga2007}. A compara\c c\~ao entre MLP NumPy e MLP Sklearn sugere que essa inicializa\c c\~ao controlada, combinada ao gradiente descendente em lote, contribuiu para a converg\^encia mais est\'avel do MLP NumPy nas configura\c c\~oes mais longas, enquanto o MLP Sklearn com SGD com momento apresentou maior instabilidade.

\subsection{SGD com Momento versus Gradiente Descendente em Lote}
\label{sec:discussao_adam}

A compara\c c\~ao entre as implementa\c c\~oes NumPy (gradiente descendente em lote sem momento) e scikit-learn (SGD com momento para o MLP), concebida primariamente para verificar a corretude das implementa\c c\~oes, constitui tamb\'em um experimento natural sobre os efeitos do otimizador no desempenho em datasets pequenos. Os resultados apontam sistematicamente para uma vantagem do gradiente descendente em lote nesse regime: em todos os casos em que o MLP Sklearn degradou (Banco 1 C3 e Banco 2 C2--C3), o MLP NumPy ou manteve o desempenho ou melhorou.

{\sloppy O mecanismo por tr\'as dessa diferen\c ca est\'a bem documentado na literatura \cite{silva2016,dsa2022}: o SGD com momento acumula uma velocidade que combina o gradiente do mini-lote com a dire\c c\~ao de atualiza\c c\~ao anterior, o que suaviza a converg\^encia mas tamb\'em significa que o otimizador continua impulsionando os par\^ametros mesmo quando o gradiente instant\^aneo se torna pequeno e o modelo j\'a est\'a em regime de memoriza\c c\~ao dos dados de treino. O gradiente descendente em lote, por calcular o gradiente exato sobre todos os exemplos de treinamento a cada \'epoca sem momento acumulado, n\~ao possui esse efeito de persist\^encia: quando a fun\c c\~ao de custo de treinamento se estabiliza, as atualiza\c c\~oes naturalmente diminuem, funcionando como uma forma impl\'icita de controle do \textit{overfitting}.\par}

Essa observa\c c\~ao tem relev\^ancia metodol\'ogica: na escolha entre otimizadores para problemas de classifica\c c\~ao cl\'inica com dados tabulares de pequeno a m\'edio porte, o gradiente descendente em lote sem momento pode ser prefer\'ivel ao SGD com momento, n\~ao por ser te\'oricamente superior em larga escala, mas por ser mais conservador em regimes de baixa quantidade de dados.

\subsection{O Papel da Binariza\c c\~ao: Evid\^encia Emp\'irica do Experimento Complementar}
\label{sec:discussao_binarizacao}

A hip\'otese de que a binariza\c c\~ao total teria beneficiado o ADALINE em detrimento do MLP foi testada empiricamente no experimento complementar da Se\c c\~ao~\ref{sec:resultados_bruto}. Os resultados revelaram que a binariza\c c\~ao prejudicou o ADALINE mais do que o beneficiou: no Banco~2, o ADALINE NumPy com dados brutos atingiu 82,22\% de acur\'acia, contra uma faixa de 68,89\% a 75,56\% com dados binarizados, ganho de at\'e 13,3 pontos percentuais. Mesmo com representa\c c\~ao cont\'inua, no entanto, o MLP Sklearn colapsou para 45,56\% de acur\'acia no Banco~2 em C3, confirmando que o problema central n\~ao \'e o tipo de representa\c c\~ao, mas o volume insuficiente de dados para regularizar um otimizador com momento acumulado.

A constata\c c\~ao central \'e que o fator determinante do desempenho relativo entre ADALINE e MLP n\~ao \'e a binariza\c c\~ao, mas o volume de dados dispon\'ivel. Com dados cont\'inuos, o ADALINE melhora de forma consistente em ambos os bancos; o MLP tamb\'em melhora no banco de maior volume, mas colapsa no menor independentemente da representa\c c\~ao das vari\'aveis. O volume de dados emerge como a vari\'avel latente que governa a competi\c c\~ao entre capacidade e generaliza\c c\~ao, resultado coerente com o dilema vi\'es-vari\^ancia \cite{braga2007,silva2016} e com a observa\c c\~ao de \citeonline{mesquita2020} sobre a primazia do volume de dados sobre a sofistica\c c\~ao do modelo.

\section{Di\'alogo com os Trabalhos Relacionados}
\label{sec:dialogo_literatura}

A discuss\~ao dos resultados ganha profundidade quando confrontada com os tr\^es trabalhos relacionados apresentados no Cap\'itulo~\ref{cap:trabalhos}, cada um contribuindo com uma perspectiva distinta sobre a rela\c c\~ao entre complexidade do modelo e desempenho preditivo em cardiologia.

\citeonline{ravani2024} obteve 89\% de acur\'acia com K-Vizinhos Mais Pr\'oximos (KNN) no mesmo banco de dados UCI Heart Disease utilizado aqui como Banco 1. Comparando com o melhor resultado obtido neste trabalho, 88,04\% (MLP Sklearn C2), a diferen\c ca \'e de apenas 0,96 ponto percentual. Essa proximidade \'e relevante por duas raz\~oes: primeiro, o KNN \'e um modelo completamente diferente das redes neurais, sem fase de treinamento expl\'icita e baseado em dist\^ancia no espa\c co de caracter\'isticas, o que sugere que m\'ultiplas fam\'ilias de algoritmos convergem para n\'iveis similares de desempenho nesse banco; segundo, o pr\'e-processamento adotado por Ravani difere do adotado neste trabalho, pois a binariza\c c\~ao total utilizada aqui n\~ao foi empregada l\'a, o que encontra confirma\c c\~ao emp\'irica no experimento complementar deste trabalho (Se\c c\~ao~\ref{sec:resultados_bruto}): com dados brutos, os modelos atingiram at\'e 88,77\% no Banco~1, evidenciando que a representa\c c\~ao cont\'inua de fato amplia o potencial preditivo do banco.

\citeonline{santos2023}, com arquiteturas LSTM para dados de ECG, obteve valores de precis\~ao entre 70\% e 84\% dependendo da classe cardiovascular avaliada, faixa compar\'avel aos resultados do Banco 2 deste trabalho (63\% a 79\%). Embora o dom\'inio seja diferente (ECG longitudinal versus dados tabulares cl\'inicos), o resultado \'e ilustrativo: mesmo aprendizado profundo com arquiteturas especializadas para s\'eries temporais n\~ao garantiu resultados substancialmente superiores aos obtidos com redes rasas em dados cl\'inicos de tamanho limitado.

\citeonline{mesquita2020}, em sua revis\~ao narrativa sobre intelig\^encia artificial em cardiologia, enfatizaram que a qualidade e o volume dos dados s\~ao fatores determinantes do desempenho, independentemente da sofistica\c c\~ao do modelo escolhido. Esse princ\'ipio, discutido no plano conceitual por aqueles autores, encontra respaldo emp\'irico direto nos resultados deste trabalho: a degrada\c c\~ao do MLP Sklearn no Banco 2 e a superioridade do ADALINE no regime de treinamento prolongado s\~ao consequ\^encias diretas das limita\c c\~oes de volume e representa\c c\~ao dos dados, n\~ao de defici\^encias das arquiteturas em si.

Coletivamente, os tr\^es trabalhos relacionados convergem para uma narrativa coerente: em dados cl\'inicos cardiovasculares de tamanho moderado, modelos simples s\~ao competitivos com modelos complexos, e a superioridade dos modelos de maior capacidade tende a emergir apenas quando o volume e a representa\c c\~ao das vari\'aveis justificam sua complexidade adicional.

\section{Implica\c c\~oes Cl\'inicas}
\label{sec:implicacoes_clinicas}

A an\'alise das implica\c c\~oes cl\'inicas dos resultados requer contextualizar os n\'umeros obtidos dentro de cen\'arios de uso real.

No Banco 1, o MLP NumPy C3 identificou corretamente 138 dos 153 pacientes doentes no conjunto de teste (sensibilidade de 90,20\%), deixando 15 sem diagn\'ostico. O ADALINE NumPy C3, com acur\'acia global pr\'oxima, detectou 133 pacientes (sensibilidade de 86,93\%), errando em 20. Se esses modelos fossem aplicados a uma popula\c c\~ao de 1.000 pacientes com distribui\c c\~ao similar (\~55\% doentes), a diferen\c ca de 5 FNs no conjunto de teste de 276 pacientes se escalaria para aproximadamente 18 FNs adicionais. Em termos de sa\'ude p\'ublica, essa diferen\c ca \'e clinicamente significativa: 18 pacientes card\'iacos por mil consultados n\~ao receberiam diagn\'ostico precoce, com potencial impacto em mortalidade e complica\c c\~oes cardiovasculares \cite{porto2019}.

Modelos com acur\'acia entre 85\% e 90\% em dados de triagem ambulatorial, como os do Banco 1, n\~ao s\~ao adequados como \'unica ferramenta de diagn\'ostico, mas podem ser valiosos como apoio \`a decis\~ao cl\'inica: identificar pacientes de alto risco que devem ser priorizados para exames complementares como ecocardiografia e cineangiocoronariografia, exames de maior custo e menor disponibilidade, especialmente em regi\~oes com menor infraestrutura assistencial \cite{brandao2025}. Nesse papel de triagem assistida, mesmo a sensibilidade de 90,20\% do melhor modelo deixaria em torno de 10\% dos pacientes doentes sem prioriza\c c\~ao, exigindo que o cl\'inico n\~ao baseie sua conduta exclusivamente na sa\'ida do algoritmo.

Para o Banco 2, a interpreta\c c\~ao \'e ainda mais matizada. Os resultados obtidos, sensibilidade m\'axima de 68,97\% (MLP NumPy C1) com acur\'acia de 78,89\%, refletem a dificuldade natural do problema de progn\'ostico de mortalidade: prever quais pacientes com insufici\^encia card\'iaca diagnosticada ir\~ao a \'obito com base em dados tabulares simples \'e uma tarefa com incerteza intr\'inseca elevada, pois a mortalidade depende de fatores n\~ao representados no banco de dados \cite{chicco2020,rohde2018}. Nesse contexto, a sensibilidade de 69\% significa que aproximadamente 1 em cada 3 pacientes que vieram a \'obito n\~ao seria identificado como de alto risco pelo melhor modelo, resultado inaceit\'avel como \'unica ferramenta de progn\'ostico, mas potencialmente \'util em combina\c c\~ao com escore cl\'inico e avalia\c c\~ao especializada.

Cabe ressaltar, ainda, que a vari\'avel \texttt{time}, n\'umero de dias de acompanhamento cl\'inico, presente no Banco 2 constitui um caso de potencial \textit{data leakage}: pacientes que foram a \'obito tendem a apresentar tempos de acompanhamento sistematicamente menores, pois o per\'iodo de seguimento se encerra com o desfecho, tornando essa vari\'avel parcialmente determinada pelo pr\'oprio alvo de predi\c c\~ao. Essa quest\~ao, identificada e documentada no Cap\'itulo~\ref{cap:metodologia}, implica que as m\'etricas obtidas no Banco 2 podem estar artificialmente elevadas: os modelos aprendem, em parte, uma associa\c c\~ao temporal espec\'ifica do conjunto de dados, e n\~ao apenas padr\~oes cl\'inicos generalizados. Isso refor\c ca a cautela ao interpretar os resultados do Banco 2 e ao compar\'a-los com os do Banco 1 \cite{chicco2020}.

\section{Sobre a Implementa\c c\~ao Manual e a Verifica\c c\~ao Emp\'irica da Teoria}
\label{sec:implementacao_manual}

A correspond\^encia entre as implementa\c c\~oes manuais (NumPy) e as implementa\c c\~oes com biblioteca (scikit-learn) constitui evid\^encia emp\'irica da corretude das implementa\c c\~oes pr\'oprias: ADALINE NumPy e ADALINE Sklearn produziram acur\'acias id\^enticas em C1 em ambos os bancos (85,87\% no Banco~1 e 68,89\% no Banco~2) e tamb\'em em C2 no Banco~1 (86,23\%), confirmando que a regra delta, o gradiente descendente e a fun\c c\~ao sigmoide foram implementados corretamente. As fun\c c\~oes de custo distintas (MSE no NumPy e entropia cruzada no Sklearn) explicam por que a correspond\^encia de acur\'acia pode n\~ao se estender \`a distribui\c c\~ao completa de VP, VN, FP e FN.

As diferen\c cas entre MLP NumPy e MLP Sklearn decorrem de diferen\c cas leg\'itimas de otimizador (gradiente em lote vs.\ SGD com momento $\beta = 0{,}9$) e inicializa\c c\~ao, n\~ao de erros \cite{pedregosa2011}. A correspond\^encia n\~ao \'e de resultados id\^enticos, mas de comportamentos coerentes: ambos os MLPs melhoram no Banco~1 com configura\c c\~oes moderadas e ambos degradam no Banco~2 com excesso de treinamento, confirmando que as implementa\c c\~oes capturam os mesmos fen\^omenos por caminhos distintos. Esse exerc\'icio de implementa\c c\~ao do zero e verifica\c c\~ao pela correspond\^encia com a biblioteca de refer\^encia preenche uma lacuna que as bibliotecas de alto n\'ivel, por sua natureza abstrata, n\~ao oferecem: o entendimento do funcionamento interno dos modelos, n\~ao apenas de seus resultados finais.

\section{Custo Computacional do Treinamento}
\label{sec:custo_computacional}

O tempo de ajuste de pesos foi medido em execu\c c\~oes auxiliares com avalia\c c\~ao por
\'epoca desativada, de forma a isolar o custo estritamente computacional do treinamento,
sem inclus\~ao das chamadas de avalia\c c\~ao utilizadas para gerar as curvas de aprendizado.
As Tabelas~\ref{tab:tempo_banco1} e~\ref{tab:tempo_banco2} apresentam os tempos obtidos
para os quatro modelos nas tr\^es configura\c c\~oes de treinamento.

\begin{table}[H]
\centering
\caption{Tempo de ajuste de pesos (segundos), Banco 1 (513 amostras de treino).}
\label{tab:tempo_banco1}
\resizebox{\textwidth}{!}{%
\begin{tabular}{|l|c|c|c|}
\hline
\textbf{Modelo} & \textbf{C1 (500 ep.)} & \textbf{C2 (1.000 ep.)} & \textbf{C3 (3.000 ep.)} \\
\hline
ADALINE NumPy   & 0,017 s & 0,032 s & 0,096 s \\
ADALINE Sklearn & 0,025 s & 0,046 s & 0,138 s \\
MLP NumPy       & 0,051 s & 0,099 s & 0,303 s \\
MLP Sklearn     & 0,199 s & 0,393 s & 1,171 s \\
\hline
\end{tabular}%
}
\fonte{Elaborado pelo autor.}
\end{table}

\begin{table}[H]
\centering
\caption{Tempo de ajuste de pesos (segundos), Banco 2 (167 amostras de treino).}
\label{tab:tempo_banco2}
\resizebox{\textwidth}{!}{%
\begin{tabular}{|l|c|c|c|}
\hline
\textbf{Modelo} & \textbf{C1 (500 ep.)} & \textbf{C2 (1.000 ep.)} & \textbf{C3 (3.000 ep.)} \\
\hline
ADALINE NumPy   & 0,009 s & 0,018 s & 0,053 s \\
ADALINE Sklearn & 0,006 s & 0,011 s & 0,030 s \\
MLP NumPy       & 0,023 s & 0,046 s & 0,133 s \\
MLP Sklearn     & 0,065 s & 0,129 s & 0,389 s \\
\hline
\end{tabular}%
}
\fonte{Elaborado pelo autor.}
\end{table}

O primeiro padr\~ao observado \'e o escalonamento aproximadamente linear do tempo
com o n\'umero de \'epocas: de C1 para C3, que representa um aumento de seis vezes no
n\'umero de \'epocas (de 500 para 3.000), o tempo de cada modelo cresce em fator entre 5,5 e
6,5 vezes em todos os casos. Esse comportamento \'e esperado, pois o custo computacional
por \'epoca \'e constante para uma mesma arquitetura e conjunto de dados, e confirma que os
algoritmos n\~ao apresentam degrada\c c\~ao de desempenho ao longo das itera\c c\~oes.

O segundo padr\~ao refere-se \`a rela\c c\~ao entre ADALINE e MLP dentro de uma mesma
estrat\'egia de implementa\c c\~ao. O MLP NumPy \'e aproximadamente 3 vezes mais lento que
o ADALINE NumPy em ambos os bancos, e o MLP Sklearn \'e cerca de 4 vezes mais
lento que o ADALINE Sklearn. Essa diferen\c ca \'e estrutural: o MLP realiza dois percursos
matriciais por \'epoca (propaga\c c\~ao direta e retropropaga\c c\~ao), enquanto o ADALINE realiza
apenas um (c\'alculo direto do gradiente), exigindo o dobro de opera\c c\~oes de multiplica\c c\~ao
matricial a cada itera\c c\~ao de treinamento \cite{braga2007}.

O terceiro padr\~ao diz respeito \`a compara\c c\~ao entre as implementa\c c\~oes NumPy e
scikit-learn para cada arquitetura. Para o ADALINE, os tempos s\~ao compar\'aveis: no
Banco 1, o ADALINE NumPy \'e levemente mais r\'apido (0,096 s versus 0,138 s em
C3); no Banco 2, a rela\c c\~ao se inverte e o ADALINE Sklearn apresenta menor tempo
(0,030 s versus 0,053 s em C3). Essa varia\c c\~ao decorre das estrat\'egias de atualiza\c c\~ao
de cada implementa\c c\~ao: o ADALINE NumPy utiliza gradiente descendente em lote,
calculando o gradiente em uma \'unica opera\c c\~ao vetorial sobre todas as 513 amostras
de treino (Banco 1) ou 167 (Banco 2) a cada \'epoca; o \texttt{SGDClassifier} do scikit-learn
realiza uma atualiza\c c\~ao de peso por amostra a cada \'epoca (SGD puro, sem mini-lotes),
percorrendo 513 atualiza\c c\~oes por \'epoca no Banco 1 e 167 no Banco 2, por meio de c\'odigo
Cython otimizado. Em conjuntos maiores, a opera\c c\~ao vetorial do NumPy se beneficia de
rotinas BLAS de alto desempenho; em conjuntos menores, o la\c co Cython por amostra
\'e competitivo com o custo de inicializa\c c\~ao da opera\c c\~ao matricial. Em termos absolutos,
nenhuma das duas implementa\c c\~oes apresentou custo relevante: mesmo o pior caso do
ADALINE (Sklearn C3, Banco 1) completou o treinamento em 0,138~s.

Para o MLP, o \texttt{MLPClassifier} do scikit-learn \'e consistentemente mais lento que
o MLP NumPy nos dois bancos, em fator de aproximadamente 4 vezes. Diferentemente
do \texttt{SGDClassifier}, que atualiza os pesos amostra por amostra, o \texttt{MLPClassifier} utiliza
mini-lotes de tamanho $\min(200, n)$: para o Banco 1 (513 amostras), isso resulta em 3
mini-lotes por \'epoca, cada um exigindo propaga\c c\~ao direta e retropropaga\c c\~ao sobre at\'e 200
amostras; para o Banco 2 (167 amostras), o mini-lote coincide com o lote completo. A
diferen\c ca de velocidade entre MLP Sklearn e MLP NumPy n\~ao se explica pelo n\'umero
de amostras processadas por \'epoca, que \'e equivalente nos dois casos, mas pelo overhead
interno do \texttt{MLPClassifier} (verifica\c c\~oes de converg\^encia, reconfigura\c c\~ao de estado entre
chamadas) e pela maior lat\^encia por mini-lote em rela\c c\~ao \`a opera\c c\~ao matricial \'unica do
NumPy.

Em todos os cen\'arios experimentados, os tempos de treinamento s\~ao inferiores a
1,2 segundos, o que torna o custo computacional uma dimens\~ao irrelevante para a escolha
entre modelos no contexto de dados tabulares de pequeno a m\'edio porte. A relev\^ancia
pr\'atica do custo computacional emergeria apenas em cen\'arios de retreinamento cont\'inuo
em larga escala ou de implanta\c c\~ao em dispositivos embarcados com recursos restritos
\cite{dsa2022}.
